\documentclass[12pt]{article}
\usepackage[utf8]{inputenc}
\usepackage[T1]{fontenc}
\usepackage{hyperref}
\title{Walman Language Resources for K-12 Curriculum Development}
\author{Compiled by Claude Sonnet 5 (claude-sonnet-5) for Andr\'as Kornai}
\date{}
\begin{document}
\maketitle

\section{Introduction}
Walman (also spelled Valman; speakers also called Koroko) is a Torricelli language of Sandaun Province, Papua New Guinea, east of Aitape. Wikipedia cites about 1,700 speakers (2003). Ethnologue-derived data cited by the Endangered Languages Project say that speakers are mostly over 60 and that most children under 20 do not understand the language. Dryer's fieldwork notes that Tok Pisin has replaced Walman for most communication in Lemieng village.

\section{Primary Educational Resources}
No primers, textbooks or school programmes in Walman were found in the searches.

\section{Linguistic and Reference Materials}
\begin{itemize}
\item Wikipedia, Walman language: \url{https://en.wikipedia.org/wiki/Walman_language}
\item Matthew Dryer, work on Papuan languages (University at Buffalo; Walman research with Lea Brown): \url{https://www.acsu.buffalo.edu/~dryer/papuan.htm}
\item Culture in Crisis, field research in Sandaun Province (an ELDP-funded project aiming at a grammar and dictionary of Walman): \url{https://cultureincrisis.org/projects/field-research-in-sandaun-province-papua-new-guinea}
\item Walman numerals page, Max Planck Institute (EVA): \url{https://lingweb.eva.mpg.de/channumerals/Walman.htm}
\end{itemize}

\section{Religious and Community Resources}
\begin{itemize}
\item Joshua Project profile, Valman/Koroko in Papua New Guinea: \url{https://joshuaproject.net/people_groups/15766}
\end{itemize}

\section{Digital and Audio Resources}
\begin{itemize}
\item Dryer's photographs from Lemieng village: \url{https://www.acsu.buffalo.edu/~dryer/walmanPics.htm}
\item Endangered Languages Project entry: \url{https://endangeredlanguages.com/lang/10550?hl=fr}
\end{itemize}
Whether the documentation project deposited audio or text corpora was not established in the searches.

\section{Curriculum Development Context}
The language shows strong shift to Tok Pisin. A grammar and dictionary may result from the documentation work, and contact with the researchers would be the natural first step; no contact details are given here. A curriculum would need new materials developed with the community.
\end{document}